\section{Dionysius in the Treatise}\label{app:ch-in-summa}

Aquinas treats Dionysius as the standing authority on the angels,
citing him at every stage of the two treatises. In \ST{I qq.~50--64} on
the angelic nature and \ST{I qq.~106--114} on the divine government of
the angels, with \ST{II-II q.~8 a.~1} besides, he cites the Dionysian
corpus explicitly ninety-eight times. By work the count is: \emph{Celestial Hierarchy},
58 (including one citation by the Latin title \emph{Hier.~Ang.}~x at
\ST{I q.~50 a.~4 s.c.}); \emph{Divine Names}, 36; \emph{Ecclesiastical
Hierarchy}, 5. (The citation at \ST{I q.~106 a.~3 s.c.} names two works
and is counted under each.) Dionysius is named without a work in
thirteen further places, most of them the extended comparison of the
two orderings of the choirs in \ST{I q.~108 a.~6}. No letter and no
chapter of the \emph{Mystical Theology} is cited in these questions.
The \emph{Celestial Hierarchy} is cited from all of its chapters except
the second; chapters 7 (15 times), 8, 6, 3, and 14 carry the most
weight. Gregory's \emph{Homiliae in Evangelia} 34, the other great
witness to the nine orders, is cited alongside at \ST{I q.~108 aa.~5--6}
and \ST{I q.~112 a.~2} (on both witnesses, \S\ref{sec:ch}; on the
orders, Appendix~\ref{app:orders}).

{\footnotesize
\begin{longtable}{@{}>{\raggedright\arraybackslash}p{0.14\linewidth}>{\raggedright\arraybackslash}p{0.10\linewidth}>{\raggedright\arraybackslash}p{0.16\linewidth}>{\raggedright\arraybackslash}p{0.52\linewidth}@{}}
\toprule
\emph{Summa} & \emph{Div.} & \emph{Dionysius} & \emph{Use} \\
\midrule
\endfirsthead
\toprule
\emph{Summa} & \emph{Div.} & \emph{Dionysius} & \emph{Use} \\
\midrule
\endhead
\bottomrule
\endfoot
I q.~50 a.~2 & s.c. & \emph{DN} 4 & First creatures immaterial and incorporeal \\
I q.~50 a.~3 & obj.~4 & \emph{DN} 4 & Intelligences subsist by the divine rays \\
I q.~50 a.~3 & co. & \emph{CH} 14 & Angelic hosts surpass material number \\
I q.~50 a.~4 & s.c. & \emph{CH} 10 (\emph{Hier.\ Ang.}~x) & First, middle, and last in one order \\
I q.~50 a.~5 & s.c. & \emph{DN} 4 & Angels free from corruption and death \\
I q.~51 a.~1 & s.c. & \emph{DN} 4 & The angels are understood to be incorporeal \\
I q.~51 a.~3 & ad~2 & \emph{CH} (15) & Members of assumed bodies signify powers \\
I q.~54 a.~2 & s.c. & \emph{DN} 4 & To understand is not a movement \\
I q.~54 a.~2 & co. & \emph{DN} 5 & God alone comprehends all things in himself \\
I q.~54 a.~3 & obj.~1 & --- & Dionysius styles angels ``intellects'' and ``minds'' \\
I q.~54 a.~3 & s.c. & \emph{CH} 11 & Substance, power, and operation divided \\
I q.~54 a.~5 & obj.~3 & \emph{DN} 4 & A ``perverted phantasy'' in the demons \\
I q.~55 a.~1 & obj.~1 & \emph{DN} 7 & Angels know earthly things by their nature \\
I q.~55 a.~1 & obj.~3 & \emph{DN} 4 & God enfolds the whole in the whole \\
I q.~55 a.~1 & s.c. & \emph{DN} 4 & Angels enlightened by the forms of things \\
I q.~55 a.~2 & s.c. & \emph{DN} 7 & Knowledge not gathered from sensible things \\
I q.~55 a.~3 & s.c. & \emph{CH} 12 & Higher angels have more universal knowledge \\
I q.~56 a.~1 & obj.~1 & \emph{CH} 6 & ``The angels do not know their own powers'' \\
I q.~56 a.~3 & obj.~1 & \emph{DN} 1 & God placed above all heavenly minds \\
I q.~56 a.~3 & obj.~3 & \emph{DN} 7 & No knowledge of God from sensible things \\
I q.~56 a.~3 & ad~1 & --- & Dionysius speaks of the knowledge of comprehension \\
I q.~57 a.~1 & co. & \emph{DN} 1; \emph{CH} 4 & Angels nearest God; things pre-exist in them \\
I q.~57 a.~5 & obj.~3 & \emph{CH} 4 & Men enlightened by the angels \\
I q.~57 a.~5 & s.c. & \emph{CH} 7 & Angels question Jesus and learn the mystery \\
I q.~57 a.~5 & ad~1 & --- & The passage of Dionysius already quoted \\
I q.~58 a.~1 & obj.~1 & \emph{DN} 4 & Angelic minds moved by understanding \\
I q.~58 a.~3 & s.c. & \emph{DN} 7 & No discourse from the common to the particular \\
I q.~58 a.~3 & co. & --- & Dionysius calls the angels ``heavenly minds'' \\
I q.~58 a.~4 & s.c. & \emph{DN} 7 & Clear simplicity of the angelic concepts \\
I q.~58 a.~4 & co. & \emph{DN} 4 & The angel a pure and most clear mirror \\
I q.~58 a.~5 & obj.~1 & \emph{DN} 4 & A ``perverted fancy'' in the demons \\
I q.~58 a.~5 & obj.~2 & \emph{EH} 6 & Nescience possible in the angels \\
I q.~58 a.~5 & obj.~3 & \emph{DN} 7 & Falsehood affirmed of the demons \\
I q.~59 a.~4 & obj.~1 & \emph{DN} 4 & Fury and concupiscence in the demons \\
I q.~60 a.~1 & obj.~1 & \emph{DN} 4 & Natural love distinguished from intellectual \\
I q.~60 a.~2 & obj.~1 & \emph{DN} 4 & Dilection identified with intellectual love \\
I q.~60 a.~3 & obj.~2 & \emph{DN} 4 & Love ``a uniting and a binding power'' \\
I q.~60 a.~3 & ad~2 & --- & ``Uniting'' and ``binding'' show love's derivation \\
I q.~63 a.~4 & s.c. & \emph{DN} 4 & The demons are not naturally wicked \\
I q.~63 a.~4 & ad~3 & \emph{DN} 4 & Not evil in the dog to be fierce \\
I q.~63 a.~7 & obj.~1 & \emph{CH} 6--7 & Cherubim under the order of Seraphim \\
I q.~64 a.~1 & s.c. & \emph{DN} 4 & Gifts bestowed on the demons unchanged \\
I q.~64 a.~1 & co. & \emph{DN} 4 & Natural knowledge undiminished in them \\
I q.~64 a.~2 & obj.~5 & \emph{DN} 4 & Demons desire to be, to live, to understand \\
I q.~106 a.~1 & s.c. & \emph{CH} 8 & Second hierarchy enlightened by the first \\
I q.~106 a.~1 & co. & \emph{CH} 7; \emph{CH} 15 & Taught by higher minds; knowledge divided for the lower \\
I q.~106 a.~1 & ad~1 & \emph{DN} 4 & Angels enlightened by the types of things \\
I q.~106 a.~2 & obj.~1 & --- & As one angel enlightens, so he cleanses and perfects \\
I q.~106 a.~2 & obj.~2 & \emph{CH} 7 & The names designate the properties; Seraphim kindle \\
I q.~106 a.~2 & ad~1 & \emph{EH} 6 & Chastening of the inferior is an enlightening \\
I q.~106 a.~3 & s.c. & \emph{CH} 4; \emph{EH} 5 & Inferior things led to God by the superior \\
I q.~106 a.~4 & obj.~1 & \emph{CH} 12 & Higher knowledge more universal than the lower \\
I q.~106 a.~4 & obj.~2 & \emph{CH} 7 & ``Who is this King of glory?'' \\
I q.~106 a.~4 & s.c. & \emph{CH} 15 & The superior communicates the gift received \\
I q.~107 a.~2 & s.c. & \emph{CH} 7 & Lower angels ask the higher, ``Who is this King of Glory?'' \\
I q.~107 a.~5 & obj.~3 & \emph{CH} 15 & Each communicates what he learns to the others \\
I q.~108 a.~1 & obj.~2 & \emph{CH} 3 & Hierarchy is order, knowledge, and action \\
I q.~108 a.~1 & s.c. & \emph{CH} 6 & Three hierarchies of angels distinguished \\
I q.~108 a.~1 & co. & \emph{CH} 1; \emph{CH} 7; \emph{CH} 3 & Sensible signs; the vestibule of God; cleansing, enlightening, perfecting \\
I q.~108 a.~2 & obj.~1 & \emph{CH} 3 & Hierarchy is order \\
I q.~108 a.~2 & obj.~3 & \emph{EH} 5 & The offices of cleansing, enlightening, perfecting \\
I q.~108 a.~2 & co. & \emph{CH} 6 & Three orders in every hierarchy \\
I q.~108 a.~3 & co. & \emph{CH} 6 & Our knowledge of the angels is imperfect \\
I q.~108 a.~3 & ad~1 & \emph{CH} 10 & First, middle, and last in one order \\
I q.~108 a.~4 & obj.~1 & \emph{CH} 3 & Hierarchy ``approaches a resemblance to God'' \\
I q.~108 a.~4 & obj.~2 & \emph{CH} 7 & Seraphim called ``burning,'' a mark of charity \\
I q.~108 a.~5 & co. & \emph{CH} 7 (twice) & Each order's name expresses its property \\
I q.~108 a.~5 & ad~1 & \emph{CH} 5; \emph{CH} 8 & Common name proper to the lowest; ``virtues'' a virile strength \\
I q.~108 a.~5 & ad~2 & \emph{DN} 12; \emph{CH} 8 & Dominion attributed to God; Domination a liberty from servitude \\
I q.~108 a.~5 & ad~3 & \emph{CH} 8; \emph{CH} 9 & Power an ordination; Principality a sacred leading \\
I q.~108 a.~5 & ad~4 & \emph{CH} 9 & Archangels between Principalities and Angels \\
I q.~108 a.~5 & ad~5 & \emph{CH} 7 (twice) & Seraphim from fire; Cherubim fulness of knowledge \\
I q.~108 a.~5 & ad~6 & \emph{CH} 7 & Thrones named from material seats \\
I q.~108 a.~6 & obj.~4 & \emph{CH} 9 & Principalities immediately above Archangels \\
I q.~108 a.~6 & s.c. & \emph{CH} 7 & Dionysius' placement of the three hierarchies \\
I q.~108 a.~6 & co. & --- & The orderings of Dionysius and Gregory compared \\
I q.~108 a.~6 & ad~4 & --- & Little or no real difference between them \\
I q.~109 a.~1 & co. & \emph{DN} 4 & The fallen keep their natural gifts \\
I q.~110 a.~3 & co. & \emph{DN} 7 & Ends of the first joined to principles of the second \\
I q.~111 a.~1 & obj.~1 & \emph{EH} 3 & Enlightenment attributed to baptism \\
I q.~111 a.~1 & s.c. & \emph{CH} 4 & Revelation reaches men through angelic ministry \\
I q.~111 a.~1 & co. & \emph{CH} 1 & The divine ray shrouded in sacred veils \\
I q.~112 a.~2 & s.c. & \emph{CH} 13 & Higher ranks fulfil no exterior service \\
I q.~112 a.~2 & co. & \emph{CH} 13 & Superior angels never sent to outward ministry \\
I q.~112 a.~2 & ad~2 & \emph{CH} 13 & The angel sent called ``Seraph'' equivocally \\
I q.~112 a.~4 & obj.~1 & \emph{CH} 8 & Second hierarchy enlightened by the first \\
I q.~112 a.~4 & s.c. & \emph{CH} 8 & Dominations above all subjection \\
I q.~112 a.~4 & co. & \emph{CH} 7 & Properties manifested by the names \\
I q.~112 a.~4 & ad~2 & \emph{CH} 14 & Angelic multitude surpasses material things \\
I q.~113 a.~2 & obj.~2 & \emph{CH} 4, 13 & The lower are led to God through the higher \\
I q.~113 a.~2 & obj.~3 & \emph{CH} 10 & Of all the angels one is greater than another \\
I q.~113 a.~3 & s.c. & \emph{CH} 5, 9 & Guardians belong to the lowest order \\
I q.~113 a.~3 & ad~1 & \emph{CH} 10 & First, middle, and last in each order \\
I q.~114 a.~3 & obj.~1 & \emph{DN} 4 & The demons a cause of evils \\
I q.~114 a.~3 & co. & --- & The words of Damascene and Dionysius \\
II-II q.~8 a.~1 & obj.~2 & \emph{DN} 4, 7 & Knowledge given according to capacity and mode \\
\end{longtable}}

The works divide the labor between them. In the treatise on the
angelic nature the \emph{Divine Names} predominates, because its fourth chapter on the
good and its seventh on the divine wisdom carry the doctrine of
participation; in the treatise on government the \emph{Celestial
Hierarchy} is nearly the only Dionysian source, and within it the
chapters on the triads and the names. The citations are densest at
\ST{I q.~108}: more than a quarter of them stand in the single question
on the hierarchies and orders, where Dionysius is not
merely quoted but followed as the master of the whole arrangement
(\S\ref{sec:q108}).
